\section{A joint Taylor-packing bound}\label{sec:direct-packing}

We first record an elementary estimate that accounts simultaneously
for auxiliary indices and analytic Taylor indices.

\begin{lemma}[Packing with an extra coordinate]\label{lem:direct-packing}
Let \(c(x)=\sum_{i=1}^r c_i x_i\) on \(\N^r\), where every
\(c_i>0\), and let \(S\subset\N^r\) have finite cardinality \(M\).
For \(L\geq0\), put
\[
 N_c(L)=\#\{(n,x)\in\N\times\N^r:n+c(x)\leq L\}.
\]
Then
\[
                     \sum_{x\in S}c(x)\geq LM-N_c(L).
\]
\end{lemma}
\begin{proof}
Discard the negative summands in \(\sum_{x\in S}(L-c(x))\).
For each remaining \(x\), the integers \(0\leq n\leq L-c(x)\)
give \(\lfloor L-c(x)\rfloor+1\geq L-c(x)\) augmented points.
These sets are disjoint for distinct \(x\), and their union is
contained in the set counted by \(N_c(L)\).
\end{proof}

\subsection{Uniform Taylor estimates}

On the complex polydisc \(|z|,|z'|\leq R_0=100K\), every
common entire function \(f_{a,P}\) associated to a monomial column
has modulus at most
\[
 \mathcal B_H=
  \exp\left(H\left[\frac{100K}{w_0}
                           +\frac{\log(600K)}{w_*}\right]\right).
\]
Indeed, \(1<\theta<2\), \(h+k\leq H/w_0\),
\(|\alpha|\leq H/w_*\), and
\(\binom{\alpha}{a}\leq2^{|\alpha|}\).
Cauchy's estimate gives the same bound for its coefficient vectors
in the normalized Taylor monomials
\((z/R_0)^d(z'/R_0)^e\).

On the local contours \(|t|=|s|=1/2\), the translated variables
satisfy \(|z|,|z'|\leq R_0/2\). A coefficient test
\([t^{\ell'_1}s^{\ell'_2}]\) of a normalized Taylor monomial is
therefore bounded by
\[
                         2^{\ell'_1+\ell'_2}2^{-d-e}.
\]
Expand a determinant of common coefficient tests row by row in
these absolutely convergent Taylor series. If two rows use the
same combined index \((a,d,e)\), their coefficient vectors are
proportional, and the corresponding determinant term vanishes.
Thus every nonzero term has distinct combined indices.

\subsection{The analytic upper bound}

Assign to a combined index the cost
\[
                         c(a,d,e)=\nu wa+\frac{d+e}{4},
\]
and define
\[
 N(L)=\#\left\{(n,d,e,a)\in\N^{m+3}:
                  n+\frac{d+e}{4}+\nu wa\leq L\right\}.
\]
For any fixed \(\lambda>0\), Lemma~\ref{lem:direct-packing}
gives, in every nonzero Taylor term,
\begin{equation}\label{eq:packing}
 \nu\sum_{\rm rows}wa+\frac14\sum_{\rm rows}(d+e)
                               \geq\lambda HM-N(\lambda H).
\end{equation}
Since \(\log2\geq1/2\),
\[
                  2^{-d-e}\leq e^{-(d+e)/4}e^{-(d+e)/4}.
\]
Combine the first exponential with the approximation cost in
\eqref{eq:packing}. The second exponential can be summed over all
Taylor pairs, at a cost
\((1-e^{-1/4})^{-2M}\) for the \(M\) rows.
The coefficient determinant is bounded by \(M!\mathcal B_H^M\).
The local test orders contribute at most \(2^{MH/v_0}\), and
Lemma~\ref{lem:translation} gives a combined row-choice and
scalar overhead of \(Q_H^M\).
It follows that the original nonzero determinant satisfies
\begin{equation}\label{eq:direct-analytic}
 \frac{\log|\Delta_H|}{MH}
 \leq\nu\bar b_H-\lambda+\frac{N(\lambda H)}{MH}
                  +E_{\rm tr}+E_{\rm hol}+\rho_H,
\end{equation}
where
\begin{align}
 E_{\rm hol}
   &=\frac{100K}{w_0}+\frac{\log2}{v_0}
                              +\frac{\log(600K)}{w_*},
       \label{eq:hol-error}\\
 \rho_H
   &=\frac{\log M+\log Q_H-2\log(1-e^{-1/4})}{H}
                               \longrightarrow0.
       \label{eq:overhead}
\end{align}
Here \(M!\leq M^M\), while \(\log M\) and \(\log Q_H\) are
\(O(\log H)\). The same geometric majorant justifies the infinite
Taylor expansions and their determinant sums. Every estimate is
uniform in the selected monomial columns and translated row choices.

\subsection{Comparison with the arithmetic bound}

The weighted-simplex count in dimension \(m+3\) yields
\[
 N(\lambda H)\sim
 \frac{16\lambda^{m+3}H^{m+3}}
              {(m+3)!\nu^m\prod_iw_i}.
\]
Together with \eqref{eq:counts}, this gives
\begin{equation}\label{eq:packing-limit}
 \frac{N(\lambda H)}{MH}\longrightarrow
       \frac{16v_0^2\lambda^{m+3}}{(m+3)K(\nu v)^m}.
\end{equation}
Write
\[
                            E=E_{\rm ar}+E_{\rm tr}+E_{\rm hol}.
\]
Comparing \eqref{eq:arith} and \eqref{eq:direct-analytic},
using \(\nu>1\) and the uniform bound \(\bar b_H\leq v\),
and then letting \(H\) tend to infinity through its allowed
multiples, gives the necessary inequality
\begin{equation}\label{eq:necessary}
 \lambda-1-(\nu-1)v
 \leq \frac{16v_0^2\lambda^{m+3}}{(m+3)K(\nu v)^m}+E.
\end{equation}
Define \(\lambda_*>0\) by
\begin{equation}\label{eq:lambda}
                   \lambda_*^{m+2}=\frac{K(\nu v)^m}{16v_0^2}.
\end{equation}
At this value the fraction in \eqref{eq:necessary} equals
\(\lambda_* /(m+3)\). Thus
\begin{equation}\label{eq:contradiction}
 g_m:=\frac{m+2}{m+3}\lambda_*-1-(\nu-1)v\leq E.
\end{equation}
We will choose the data so that the strict reverse inequality holds.
